% appendix_model_agent.tex — Methodology appendix
% Uses macros from quantitative_stats.tex

%%%%%%%%%%%%%%%%%%%%%%%%%%%%%%%%%%%%%%%%%%%%%%%%%%%%%%%%%%%%%%%%%%%%%%%%%%%%%%%
\subsection{Benchmark Predictability: BenchPress Methodology}
\label{app:benchpress}

We adapt the BenchPress methodology~\cite{benchpress2025} to measure which
benchmarks carry unique evaluative signal versus which are redundant (i.e.,
fully reconstructable from other benchmarks).  The original BenchPress
operates on an $83 \times 49$ model--benchmark matrix; we adapt it to our
$16 \times 50$ agent+model--benchmark matrix by preserving the core algorithmic
components while adjusting for lower sample count.

\subsubsection{LogitBenchReg}
\label{app:logit-benchreg}

For percentage-bounded scores $s \in [0,1]$, we first apply the logit
transform to linearize ceiling/floor effects:
\begin{equation}
  z = \operatorname{logit}(s) = \log\!\frac{\,\mathrm{clip}(s,\,\epsilon,\,1{-}\epsilon)\,}
       {1 - \mathrm{clip}(s,\,\epsilon,\,1{-}\epsilon)},
  \quad \epsilon = 0.005.
  \label{eq:logit-transform}
\end{equation}

For each target benchmark~$b_j$, we identify the $k{=}5$ predictor benchmarks
$\{b_{j_1}, \ldots, b_{j_5}\}$ that maximize the training-set $R^2$ when
predicting $b_j$ via univariate linear regression in logit space:
\begin{equation}
  R^2(b_{j'} \to b_j)
  = 1 - \frac{\sum_i (z_{i,j} - \hat{z}_{i,j})^2}{\sum_i (z_{i,j} - \bar{z}_j)^2},
  \quad
  \hat{z}_{i,j} = \alpha_{j'} \cdot z_{i,j'} + \beta_{j'}.
  \label{eq:benchreg-r2}
\end{equation}
Only predictors with $R^2 \ge 0.1$ are retained.  For a held-out system~$i$
with known scores on some subset of these predictors, the final prediction is
an $R^2$-weighted average:
\begin{equation}
  \hat{z}_{i,j}^{\mathrm{BenchReg}}
  = \frac{\sum_{j' \in \mathcal{K}_i}\, R^2(j' \!\to\! j) \cdot (\alpha_{j'}\, z_{i,j'} + \beta_{j'})}
         {\sum_{j' \in \mathcal{K}_i}\, R^2(j' \!\to\! j)},
  \label{eq:benchreg-pred}
\end{equation}
where $\mathcal{K}_i \subseteq \{j_1,\ldots,j_5\}$ is the subset of top-$k$
predictors for which system~$i$ has observed scores.  The prediction is mapped
back to $[0,1]$ via the inverse logit $\hat{s} = \sigma(\hat{z})$.

\subsubsection{SVD-Logit (Rank-2 Soft-Impute)}
\label{app:svd-logit}

As a global complement, we perform iterative SVD completion in logit-then-$z$-score
space.  After logit-transforming all percentage-bounded columns, we column-wise
standardize to $z$-scores, initialize missing entries to 0, then iterate:
\begin{enumerate}
  \item Compute truncated SVD: $\mathbf{M} \approx \mathbf{U}_r\,\mathbf{\Sigma}_r\,\mathbf{V}_r^\top$ with $r{=}2$.
  \item Replace only the missing entries with the rank-$r$ approximation;
        observed entries remain unchanged.
  \item Repeat until relative change $< 10^{-4}$.
\end{enumerate}
After convergence, predictions are de-standardized and inverse-logit-transformed
back to $[0,1]$.

The rank $r{=}2$ is chosen following the original BenchPress finding that rank~2
is optimal for sparse LLM benchmark matrices---higher ranks overfit.

\subsubsection{BenchPress Blend}
\label{app:benchpress-blend}

The final predictor combines both methods:
\begin{equation}
  \hat{s}_{i,j}^{\mathrm{BenchPress}}
  = \alpha \cdot \hat{s}_{i,j}^{\mathrm{BenchReg}}
  + (1{-}\alpha) \cdot \hat{s}_{i,j}^{\mathrm{SVD}},
  \quad \alpha = 0.6.
  \label{eq:benchpress-blend}
\end{equation}
When LogitBenchReg has no prediction for a cell (insufficient predictor
coverage), the method falls back to SVD-Logit alone.

\subsubsection{Evaluation Protocol}
\label{app:benchpress-eval}

We measure per-benchmark predictability using a per-system holdout protocol:
\begin{enumerate}
  \item For each of 3 random folds: for every system~$i$, randomly hide 50\%
        of its observed scores.
  \item Train BenchPress on the remaining observed matrix.
  \item Predict hidden cells and record absolute errors per benchmark.
\end{enumerate}
The primary metric is the \textbf{median absolute error (MedAE)} on held-out
cells, aggregated per benchmark.  Benchmarks with higher MedAE carry more
unique signal---their scores cannot be reconstructed from other benchmarks.

\paragraph{Adaptation note.}
The original BenchPress uses 83 models as training rows; our matrix has 16
agent+model combinations.  With $k{=}5$ predictors per target, each univariate
regression is fit on $\sim$14--16 data points, which is statistically marginal
but stabilized by the $R^2$ weighting and the SVD-Logit fallback.  We report
MedAE rather than percentage error because our scores are on $[0,1]$ rather
than $[0,100]$.

%%%%%%%%%%%%%%%%%%%%%%%%%%%%%%%%%%%%%%%%%%%%%%%%%%%%%%%%%%%%%%%%%%%%%%%%%%%%%%%

\subsection{Model vs.\ Agent Effect: Methodology Details}
\label{app:model-vs-agent}

This appendix formalizes the within-family and adjusted-effect analyses
summarized in Section~\ref{sec:agent-harness}.  Let $s(m, a, b)$ denote the
benchmark score for model~$m$, agent scaffold~$a$, and benchmark~$b$.

%%%%%%%%%%%%%%%%%%%%%%%%%%%%%%%%%%%%%%%%%%%%%%%%%%%%%%%%%%%%%%%%%%%%%%%%%%%%%%%
\subsubsection{Within-Family Comparison}
\label{app:within-family}

For each model provider family $f$ (e.g., Anthropic, OpenAI, Google), let
$M_f$ be the set of models belonging to $f$ and $A_f$ the set of non-baseline
agents evaluated on that family.  The baseline agent is Terminus~2.

\paragraph{Model effect (fixing agent).}
For each benchmark~$b$, we fix the agent to Terminus~2 and measure the score
range across models within the family:
\begin{equation}
  \Delta_{\mathrm{model}}^{(f,b)}
  \;=\;
  \max_{m \in M_f}\, s(m,\,\text{terminus-2},\,b)
  \;-\;
  \min_{m \in M_f}\, s(m,\,\text{terminus-2},\,b).
  \label{eq:model-effect}
\end{equation}
This requires at least two models in~$f$ with observed Terminus~2 scores on~$b$.

\paragraph{Agent effect (fixing model).}
For each model $m \in M_f$ evaluated under both Terminus~2 and another agent
$a \in A_f$ on benchmark~$b$, we compute the absolute score change:
\begin{equation}
  \delta_{\mathrm{agent}}^{(m,a,b)}
  \;=\;
  \bigl|\, s(m,\,a,\,b) - s(m,\,\text{terminus-2},\,b) \,\bigr|.
  \label{eq:agent-delta}
\end{equation}
The family-level agent effect on~$b$ averages over all available model--agent pairs:
\begin{equation}
  \Delta_{\mathrm{agent}}^{(f,b)}
  \;=\;
  \frac{1}{|\mathcal{P}_{f,b}|}
  \sum_{(m,a)\in \mathcal{P}_{f,b}}
  \delta_{\mathrm{agent}}^{(m,a,b)},
  \label{eq:agent-effect}
\end{equation}
where $\mathcal{P}_{f,b} = \{(m,a) : m \in M_f,\; a \in A_f,\; \text{both }
s(m,a,b) \text{ and } s(m,\text{terminus-2},b) \text{ are observed}\}$.

\paragraph{Dominance classification.}
For a given $(f, b)$ pair, the benchmark is classified as
\emph{model-dominant} if $\Delta_{\mathrm{model}}^{(f,b)} >
\Delta_{\mathrm{agent}}^{(f,b)}$, and \emph{agent-dominant} otherwise.

\paragraph{Family-level summary.}
We aggregate across benchmarks within each family using the mean:
\begin{align}
  \overline{\Delta}_{\mathrm{model}}^{(f)}
  &= \frac{1}{|B_f|}\sum_{b \in B_f} \Delta_{\mathrm{model}}^{(f,b)}, &
  \overline{\Delta}_{\mathrm{agent}}^{(f)}
  &= \frac{1}{|B_f|}\sum_{b \in B_f} \Delta_{\mathrm{agent}}^{(f,b)},
  \label{eq:family-summary}
\end{align}
where $B_f$ is the set of benchmarks with both model-effect and agent-effect
estimates for family~$f$.  The model-to-agent ratio is
$\overline{\Delta}_{\mathrm{model}}^{(f)} / \overline{\Delta}_{\mathrm{agent}}^{(f)}$,
and the model-dominant percentage is the fraction of benchmarks in~$B_f$
classified as model-dominant.

Table~\ref{tab:within-family-summary} reports these statistics for the three
multi-model families.

\begin{table}[t]
  \centering
  \caption{Within-family model vs.\ agent effect comparison.  The model effect
    is the score range across models within the family on Terminus~2; the agent
    effect is the mean absolute score change from switching the agent while
    holding the model fixed.  Ratio and model-dominant percentage are computed
    per benchmark, then averaged.}
  \label{tab:within-family-summary}
  \begin{tabular}{lcccccc}
    \toprule
    Family
      & Models
      & Agents
      & $\overline{\Delta}_{\mathrm{model}}$
      & $\overline{\Delta}_{\mathrm{agent}}$
      & Ratio
      & Model-dom.\ \% \\
    \midrule
    Anthropic
      & 3
      & claude-code
      & \AnthropicModelEffectMean{}
      & \AnthropicAgentEffectMean{}
      & \AnthropicModelAgentRatio{}$\times$
      & \AnthropicModelWinPct\% \\
    OpenAI
      & 3
      & codex
      & \OpenAIModelEffectMean{}
      & \OpenAIAgentEffectMean{}
      & \OpenAIModelAgentRatio{}$\times$
      & \OpenAIModelWinPct\% \\
    Google
      & 2
      & gemini-cli
      & \GoogleModelEffectMean{}
      & \GoogleAgentEffectMean{}
      & \GoogleModelAgentRatio{}$\times$
      & \GoogleModelWinPct\% \\
    \bottomrule
  \end{tabular}
\end{table}

Table~\ref{tab:within-family-example} illustrates the per-benchmark comparison
for one representative family.

\begin{table}[t]
  \centering
  \caption{Example within-family comparison for the Anthropic family
    (claude-opus-4-6, claude-sonnet-4-6, claude-haiku-4-5).
    \emph{Model effect}: score range across the three models on Terminus~2.
    \emph{Agent effect}: mean $|\Delta|$ from switching to Claude~Code while
    holding each model fixed.  Only a representative subset of benchmarks is
    shown.}
  \label{tab:within-family-example}
  \small
  \begin{tabular}{lcccc}
    \toprule
    & \multicolumn{2}{c}{Model effect (Terminus-2)}
    & \multicolumn{2}{c}{Agent effect (fix model)} \\
    \cmidrule(lr){2-3}\cmidrule(lr){4-5}
    Benchmark & Score range & Rank & Mean $|\Delta|$ & Rank \\
    \midrule
    \multicolumn{5}{c}{\emph{Auto-populated from pipeline output ---
      see \texttt{benchmark\_within\_family\_model\_vs\_agent.csv}}} \\
    \bottomrule
  \end{tabular}
\end{table}


%%%%%%%%%%%%%%%%%%%%%%%%%%%%%%%%%%%%%%%%%%%%%%%%%%%%%%%%%%%%%%%%%%%%%%%%%%%%%%%
\subsubsection{Adjusted Group Effects}
\label{app:adjusted-effects}

The within-family analysis is restricted to families with multiple models.  To
obtain a single ranking across all \NumModels{} models and \NumAgents{} agents,
we use an adjusted group effect that controls for confounders.

Let $\mathbf{y}$ be the vector of benchmark-relative (robustly normalized)
scores, one entry per observed (model,\,agent,\,benchmark) triple.  We fit a
linear model using the confounding factors as one-hot-encoded predictors:
\begin{equation}
  \hat{\mathbf{y}} = \mathbf{X}_{\text{controls}}\,\boldsymbol{\beta},
  \label{eq:confound-fit}
\end{equation}
where $\mathbf{X}_{\text{controls}}$ contains dummy variables for the
confounding dimensions (with one level dropped per factor).

\paragraph{Model effects.}
To estimate adjusted model effects, the controls are \emph{agent} and
\emph{benchmark}:
\begin{equation}
  y^{\mathrm{adj}}_i
  = y_i - \hat{y}_i + \bar{y},
  \qquad
  \hat{y}_i = [\mathbf{X}_{\text{agent, benchmark}}\,\hat{\boldsymbol{\beta}}]_i.
  \label{eq:adjusted-model}
\end{equation}
The adjusted model effect for model~$m$ is then
$\bar{y}^{\mathrm{adj}}_m = \frac{1}{n_m}\sum_{i:\,\text{model}_i = m}
y^{\mathrm{adj}}_i$, where $n_m$ is the number of observations for model~$m$.
This removes the influence of which agents and benchmarks each model was paired
with, isolating the model's contribution.

\paragraph{Agent effects.}
Symmetrically, to estimate adjusted agent effects, the controls are
\emph{model} and \emph{benchmark}:
\begin{equation}
  y^{\mathrm{adj}}_i
  = y_i - \hat{y}_i + \bar{y},
  \qquad
  \hat{y}_i = [\mathbf{X}_{\text{model, benchmark}}\,\hat{\boldsymbol{\beta}}]_i.
  \label{eq:adjusted-agent}
\end{equation}

\paragraph{Interpretation.}
The adjusted model effect span is \ModelEffectSpan{} robust-scaled units (from
\TopModel{} at $+$\TopModelEffect{} to \BottomModel{} at
\BottomModelEffect{}), compared to \AgentEffectSpan{} for agents.  The
$>$10$\times$ ratio confirms that model selection is the dominant lever even
after controlling for the unbalanced agent--benchmark evaluation matrix.

\begin{table}[t]
  \centering
  \caption{Adjusted model effects.  Each model's mean benchmark-relative score
    is adjusted by removing the confounding effect of which agents and
    benchmarks it was paired with.  Units are robust-scaled benchmark-relative
    scores (0~=~benchmark median).}
  \label{tab:adjusted-model-effects}
  \small
  \begin{tabular}{lcc}
    \toprule
    Model & Adjusted mean & Observations \\
    \midrule
    \multicolumn{3}{c}{\emph{Auto-populated from pipeline output ---
      see \texttt{benchmark\_model\_adjusted\_effects.csv}}} \\
    \bottomrule
  \end{tabular}
\end{table}

\begin{table}[t]
  \centering
  \caption{Adjusted agent effects.  Same methodology as
    Table~\ref{tab:adjusted-model-effects}, controlling for model and
    benchmark.}
  \label{tab:adjusted-agent-effects}
  \small
  \begin{tabular}{lcc}
    \toprule
    Agent & Adjusted mean & Observations \\
    \midrule
    \multicolumn{3}{c}{\emph{Auto-populated from pipeline output ---
      see \texttt{benchmark\_agent\_adjusted\_effects.csv}}} \\
    \bottomrule
  \end{tabular}
\end{table}
